% !TEX program = xelatex
% ===========================================================================
%  《主题》—— 中文数学习题课讲义（Loom 模板 · 课前自学版）
%  必须用 XeLaTeX 编译两遍；QA 用 scripts/qa_check.py。
%  第一行的魔法注释让 VS Code (LaTeX Workshop) 走 XeLaTeX——别删（默认是
%  pdfLaTeX，会报 xeCJK requires XeTeX）。
% ===========================================================================
\documentclass{loom}

% 挑战题 ★ 用 Zapf Dingbats 绘制，不依赖系统字体（Optima 等）是否含 U+2605
\usepackage{pifont}

% 页眉线（可选）；老师偏好极简时可整行注释掉
%\runningthread{《主题》 · 课前自学讲义}

% ---------------------------------------------------------------------------
%  讲义专用组件：题面（茜红结）、解析（靛蓝结）、四段式解答标签
%  —— 这些宏定义住在 main.tex（不在 cls），以便随讲义微调。
% ---------------------------------------------------------------------------
% 题目框：breakable=false 保证一题不跨页。默认不留答题空白区（老师偏好）。
% —— 若要做「打印作答版」，取消下面 \answerspace 定义与 problem 框 after upper
%    一行的注释即可，题末会追加一块纯白书写面板（默认 5 行，可 \answerspace[8]）。
% \newcommand{\answerspace}[1][5]{\par\medskip\noindent
%   \colorbox{white}{\parbox[t][#1\baselineskip][t]{\dimexpr\linewidth-2\fboxsep\relax}{\strut}}}
\newtcolorbox{problem}[1]{knotmadder, breakable=false,
  before upper={{\headingfont\bfseries\color{madder}#1}\par\smallskip}%
  % , after upper={\answerspace}   % 打印作答版：取消本行注释
}
\newtcolorbox{solution}[1]{knotindigo,
  before upper={{\headingfont\bfseries\color{indigo}#1}\par\smallskip}}

\newcommand{\think}{\par\smallskip\noindent
  {\headingfont\bfseries\color{indigo}〔怎么想〕}\enspace}
\newcommand{\solve}{\par\smallskip\noindent
  {\headingfont\bfseries\color{inkiron}〔解〕}\enspace}
\newcommand{\solvehead}[1]{\par\smallskip\noindent
  {\headingfont\bfseries\color{inkiron}〔#1〕}\enspace}
\newcommand{\checkup}{\par\smallskip\noindent
  {\headingfont\bfseries\color{selvage}〔检验〕}\enspace}
\newcommand{\insight}[1]{\par\smallskip\noindent
  {\headingfont\bfseries\color{weld!60!inkiron}〔方法点睛〕}\enspace
  {\strandfont\small #1}\par}
\newcommand{\hint}[1]{\par\smallskip\noindent
  {\strandfont\small{\raisebox{-0.8pt}{\loomtile{weld}}\hspace{4pt}}%
  {\bfseries\color{weld!60!inkiron}入手提示}\enspace{\color{inkiron!85!white}#1}}\par}
% 页边旁批：放"想一想/追问"、∑/∏ 记号说明，不打断正文
\newcommand{\sidenote}[1]{\marginnote{\raggedright\strandfont\footnotesize
  \color{selvage}#1}}
\newcommand{\starred}{{\color{weld!70!inkiron}\ding{72}}}

\setlist[enumerate]{leftmargin=2.3em, itemsep=2pt, topsep=3pt}
\setlist[itemize]{leftmargin=1.6em, itemsep=2pt, topsep=3pt}


\usetikzlibrary{angles,quotes,arrows.meta}
\newif\ifteachernotes\teachernotestrue
\newif\ifshowhints\showhintstrue
\let\originalhint\hint
\renewcommand{\hint}[1]{\ifshowhints\originalhint{#1}\fi}
\newcommand{\diagram}[1]{\par\smallskip\begin{center}#1\end{center}\smallskip}
\hypersetup{pdftitle={双等腰模型（婆罗摩笈多模型）· 单章核对稿},pdfauthor={Hence},pdfsubject={初中几何模型培优讲义}}
\renewcommand{\loomcovermotto}{从图形中发现关系，在推理中理解几何}
\renewcommand{\contentsname}{目录}
% Same coordinates as the SVG/PNG diagrams.
\newcommand{\diReview}{%
\begin{tikzpicture}[scale=1.2,line width=.7pt,every node/.style={font=\small,inner sep=2pt}]
\coordinate (A) at (-1.00000000,0.00000000);
\coordinate (B) at (0.00000000,0.00000000);
\coordinate (C) at (2.00000000,0.00000000);
\coordinate (D) at (-1.00000000,2.00000000);
\coordinate (E) at (2.00000000,1.00000000);
\draw[inkiron] (A) -- (C);
\draw[inkiron] (A) -- (D);
\draw[inkiron] (B) -- (D);
\draw[inkiron] (B) -- (E);
\draw[inkiron] (C) -- (E);
\pic[draw=selvage,angle radius=2.2mm] {right angle=D--A--B};
\pic[draw=selvage,angle radius=2.2mm] {right angle=D--B--E};
\pic[draw=selvage,angle radius=2.2mm] {right angle=B--C--E};
\draw[inkiron] ($(B)!0.50000000!(D)!3pt!90:(D)$) -- ($(B)!0.50000000!(D)!3pt!-90:(D)$);
\draw[inkiron] ($(B)!0.50000000!(E)!3pt!90:(E)$) -- ($(B)!0.50000000!(E)!3pt!-90:(E)$);
\fill[inkiron] (A) circle (1.05pt);
\node at (-1.15000000,-0.27000000) {$A$};
\fill[inkiron] (B) circle (1.05pt);
\node at (0.00000000,-0.30000000) {$B$};
\fill[inkiron] (C) circle (1.05pt);
\node at (2.10000000,-0.27000000) {$C$};
\fill[inkiron] (D) circle (1.05pt);
\node at (-1.10000000,2.26000000) {$D$};
\fill[inkiron] (E) circle (1.05pt);
\node at (2.24000000,1.05000000) {$E$};
\end{tikzpicture}}

\newcommand{\diModel}{%
\begin{tikzpicture}[scale=0.47,line width=.7pt,every node/.style={font=\small,inner sep=2pt}]
\coordinate (B) at (0.00000000,0.00000000);
\coordinate (C) at (-6.00000000,2.00000000);
\coordinate (Cp) at (-2.00000000,-6.00000000);
\coordinate (D) at (3.00000000,1.00000000);
\coordinate (Dp) at (-1.00000000,3.00000000);
\fill[selvage!10] (C) -- (B) -- (Dp) -- cycle;
\fill[madder!10] (Cp) -- (B) -- (D) -- cycle;
\draw[inkiron] (B) -- (C);
\draw[inkiron] (B) -- (Cp);
\draw[inkiron] (B) -- (D);
\draw[inkiron] (B) -- (Dp);
\draw[inkiron] (C) -- (Dp);
\draw[inkiron] (Cp) -- (D);
\pic[draw=selvage,angle radius=2.2mm] {right angle=C--B--Cp};
\pic[draw=selvage,angle radius=2.2mm] {right angle=D--B--Dp};
\draw[inkiron] ($(B)!0.50000000!(C)!3pt!90:(C)$) -- ($(B)!0.50000000!(C)!3pt!-90:(C)$);
\draw[inkiron] ($(B)!0.50000000!(Cp)!3pt!90:(Cp)$) -- ($(B)!0.50000000!(Cp)!3pt!-90:(Cp)$);
\draw[inkiron] ($(B)!0.48260747!(D)!3pt!90:(D)$) -- ($(B)!0.48260747!(D)!3pt!-90:(D)$);
\draw[inkiron] ($(B)!0.51739253!(D)!3pt!90:(D)$) -- ($(B)!0.51739253!(D)!3pt!-90:(D)$);
\draw[inkiron] ($(B)!0.48260747!(Dp)!3pt!90:(Dp)$) -- ($(B)!0.48260747!(Dp)!3pt!-90:(Dp)$);
\draw[inkiron] ($(B)!0.51739253!(Dp)!3pt!90:(Dp)$) -- ($(B)!0.51739253!(Dp)!3pt!-90:(Dp)$);
\fill[inkiron] (B) circle (1.05pt);
\node at (-0.46000000,-0.45000000) {$B$};
\fill[inkiron] (C) circle (1.05pt);
\node at (-6.45000000,2.10000000) {$C$};
\fill[inkiron] (Cp) circle (1.05pt);
\node at (-2.20000000,-6.42000000) {$C'$};
\fill[inkiron] (D) circle (1.05pt);
\node at (3.36000000,1.12000000) {$D$};
\fill[inkiron] (Dp) circle (1.05pt);
\node at (-0.68000000,3.26000000) {$D'$};
\end{tikzpicture}}

\newcommand{\diArea}{%
\begin{tikzpicture}[scale=0.47,line width=.7pt,every node/.style={font=\small,inner sep=2pt}]
\coordinate (B) at (0.00000000,0.00000000);
\coordinate (C) at (-6.00000000,2.00000000);
\coordinate (Cp) at (-2.00000000,-6.00000000);
\coordinate (D) at (3.00000000,1.00000000);
\coordinate (Dp) at (-1.00000000,3.00000000);
\fill[selvage!10] (C) -- (B) -- (Dp) -- cycle;
\fill[madder!10] (Cp) -- (B) -- (D) -- cycle;
\draw[inkiron] (B) -- (C);
\draw[inkiron] (B) -- (Cp);
\draw[inkiron] (B) -- (D);
\draw[inkiron] (B) -- (Dp);
\draw[inkiron] (C) -- (Dp);
\draw[inkiron] (Cp) -- (D);
\pic[draw=selvage,angle radius=2.2mm] {right angle=C--B--Cp};
\pic[draw=selvage,angle radius=2.2mm] {right angle=D--B--Dp};
\draw[inkiron] ($(B)!0.50000000!(C)!3pt!90:(C)$) -- ($(B)!0.50000000!(C)!3pt!-90:(C)$);
\draw[inkiron] ($(B)!0.50000000!(Cp)!3pt!90:(Cp)$) -- ($(B)!0.50000000!(Cp)!3pt!-90:(Cp)$);
\draw[inkiron] ($(B)!0.48260747!(D)!3pt!90:(D)$) -- ($(B)!0.48260747!(D)!3pt!-90:(D)$);
\draw[inkiron] ($(B)!0.51739253!(D)!3pt!90:(D)$) -- ($(B)!0.51739253!(D)!3pt!-90:(D)$);
\draw[inkiron] ($(B)!0.48260747!(Dp)!3pt!90:(Dp)$) -- ($(B)!0.48260747!(Dp)!3pt!-90:(Dp)$);
\draw[inkiron] ($(B)!0.51739253!(Dp)!3pt!90:(Dp)$) -- ($(B)!0.51739253!(Dp)!3pt!-90:(Dp)$);
\coordinate (M) at (-1.80000000,0.60000000);
\coordinate (N) at (0.60000000,1.80000000);
\draw[indigo,dashed] (Dp) -- (M);
\draw[indigo,dashed] (D) -- (N);
\draw[indigo,dashed] (B) -- (N);
\pic[draw=selvage,angle radius=2.2mm] {right angle=B--M--Dp};
\pic[draw=selvage,angle radius=2.2mm] {right angle=B--N--D};
\fill[inkiron] (B) circle (1.05pt);
\node at (-0.46000000,-0.45000000) {$B$};
\fill[inkiron] (C) circle (1.05pt);
\node at (-6.45000000,2.10000000) {$C$};
\fill[inkiron] (Cp) circle (1.05pt);
\node at (-2.20000000,-6.42000000) {$C'$};
\fill[inkiron] (D) circle (1.05pt);
\node at (3.36000000,1.12000000) {$D$};
\fill[inkiron] (Dp) circle (1.05pt);
\node at (-0.68000000,3.26000000) {$D'$};
\fill[inkiron] (M) circle (1.05pt);
\node at (-2.13000000,0.38000000) {$M$};
\fill[inkiron] (N) circle (1.05pt);
\node at (0.64000000,2.16000000) {$N$};
\end{tikzpicture}}

\newcommand{\diMidpoint}{%
\begin{tikzpicture}[scale=0.45,line width=.7pt,every node/.style={font=\small,inner sep=2pt}]
\coordinate (B) at (0.00000000,0.00000000);
\coordinate (C) at (-6.00000000,2.00000000);
\coordinate (Cp) at (-2.00000000,-6.00000000);
\coordinate (D) at (3.00000000,1.00000000);
\coordinate (Dp) at (-1.00000000,3.00000000);
\fill[selvage!10] (C) -- (B) -- (Dp) -- cycle;
\fill[madder!10] (Cp) -- (B) -- (D) -- cycle;
\draw[inkiron] (B) -- (C);
\draw[inkiron] (B) -- (Cp);
\draw[inkiron] (B) -- (D);
\draw[inkiron] (B) -- (Dp);
\draw[inkiron] (C) -- (Dp);
\draw[inkiron] (Cp) -- (D);
\pic[draw=selvage,angle radius=2.2mm] {right angle=C--B--Cp};
\pic[draw=selvage,angle radius=2.2mm] {right angle=D--B--Dp};
\draw[inkiron] ($(B)!0.50000000!(C)!3pt!90:(C)$) -- ($(B)!0.50000000!(C)!3pt!-90:(C)$);
\draw[inkiron] ($(B)!0.50000000!(Cp)!3pt!90:(Cp)$) -- ($(B)!0.50000000!(Cp)!3pt!-90:(Cp)$);
\draw[inkiron] ($(B)!0.48260747!(D)!3pt!90:(D)$) -- ($(B)!0.48260747!(D)!3pt!-90:(D)$);
\draw[inkiron] ($(B)!0.51739253!(D)!3pt!90:(D)$) -- ($(B)!0.51739253!(D)!3pt!-90:(D)$);
\draw[inkiron] ($(B)!0.48260747!(Dp)!3pt!90:(Dp)$) -- ($(B)!0.48260747!(Dp)!3pt!-90:(Dp)$);
\draw[inkiron] ($(B)!0.51739253!(Dp)!3pt!90:(Dp)$) -- ($(B)!0.51739253!(Dp)!3pt!-90:(Dp)$);
\coordinate (E) at (-3.50000000,2.50000000);
\coordinate (G) at (-7.00000000,5.00000000);
\coordinate (H) at (1.51351351,-1.08108108);
\draw[indigo,dashed] (G) -- (H);
\draw[indigo,dashed] (C) -- (G);
\draw[indigo] ($(C)!0.45685445!(E)!3pt!90:(E)$) -- ($(C)!0.45685445!(E)!3pt!-90:(E)$);
\draw[indigo] ($(C)!0.50000000!(E)!3pt!90:(E)$) -- ($(C)!0.50000000!(E)!3pt!-90:(E)$);
\draw[indigo] ($(C)!0.54314555!(E)!3pt!90:(E)$) -- ($(C)!0.54314555!(E)!3pt!-90:(E)$);
\draw[indigo] ($(E)!0.45685445!(Dp)!3pt!90:(Dp)$) -- ($(E)!0.45685445!(Dp)!3pt!-90:(Dp)$);
\draw[indigo] ($(E)!0.50000000!(Dp)!3pt!90:(Dp)$) -- ($(E)!0.50000000!(Dp)!3pt!-90:(Dp)$);
\draw[indigo] ($(E)!0.54314555!(Dp)!3pt!90:(Dp)$) -- ($(E)!0.54314555!(Dp)!3pt!-90:(Dp)$);
\draw[madder] ($(B)!0.48721276!(E)!3pt!90:(E)$) -- ($(B)!0.48721276!(E)!3pt!-90:(E)$);
\draw[madder] ($(B)!0.51278724!(E)!3pt!90:(E)$) -- ($(B)!0.51278724!(E)!3pt!-90:(E)$);
\draw[madder] ($(E)!0.48721276!(G)!3pt!90:(G)$) -- ($(E)!0.48721276!(G)!3pt!-90:(G)$);
\draw[madder] ($(E)!0.51278724!(G)!3pt!90:(G)$) -- ($(E)!0.51278724!(G)!3pt!-90:(G)$);
\fill[inkiron] (B) circle (1.05pt);
\node at (-0.46000000,-0.45000000) {$B$};
\fill[inkiron] (C) circle (1.05pt);
\node at (-6.45000000,2.10000000) {$C$};
\fill[inkiron] (Cp) circle (1.05pt);
\node at (-2.20000000,-6.42000000) {$C'$};
\fill[inkiron] (D) circle (1.05pt);
\node at (3.36000000,1.12000000) {$D$};
\fill[inkiron] (Dp) circle (1.05pt);
\node at (-0.68000000,3.26000000) {$D'$};
\fill[inkiron] (E) circle (1.05pt);
\node at (-3.45000000,2.93000000) {$E$};
\fill[inkiron] (G) circle (1.05pt);
\node at (-7.20000000,5.35000000) {$G$};
\fill[inkiron] (H) circle (1.05pt);
\node at (1.81351351,-1.31108108) {$H$};
\end{tikzpicture}}

\newcommand{\diConverse}{%
\begin{tikzpicture}[scale=0.52,line width=.7pt,every node/.style={font=\small,inner sep=2pt}]
\coordinate (B) at (0.00000000,0.00000000);
\coordinate (C) at (-6.00000000,2.00000000);
\coordinate (Cp) at (-2.00000000,-6.00000000);
\coordinate (D) at (3.00000000,1.00000000);
\coordinate (Dp) at (-1.00000000,3.00000000);
\fill[selvage!10] (C) -- (B) -- (Dp) -- cycle;
\fill[madder!10] (Cp) -- (B) -- (D) -- cycle;
\draw[inkiron] (B) -- (C);
\draw[inkiron] (B) -- (Cp);
\draw[inkiron] (B) -- (D);
\draw[inkiron] (B) -- (Dp);
\draw[inkiron] (C) -- (Dp);
\draw[inkiron] (Cp) -- (D);
\pic[draw=selvage,angle radius=2.2mm] {right angle=C--B--Cp};
\pic[draw=selvage,angle radius=2.2mm] {right angle=D--B--Dp};
\draw[inkiron] ($(B)!0.50000000!(C)!3pt!90:(C)$) -- ($(B)!0.50000000!(C)!3pt!-90:(C)$);
\draw[inkiron] ($(B)!0.50000000!(Cp)!3pt!90:(Cp)$) -- ($(B)!0.50000000!(Cp)!3pt!-90:(Cp)$);
\draw[inkiron] ($(B)!0.48260747!(D)!3pt!90:(D)$) -- ($(B)!0.48260747!(D)!3pt!-90:(D)$);
\draw[inkiron] ($(B)!0.51739253!(D)!3pt!90:(D)$) -- ($(B)!0.51739253!(D)!3pt!-90:(D)$);
\draw[inkiron] ($(B)!0.48260747!(Dp)!3pt!90:(Dp)$) -- ($(B)!0.48260747!(Dp)!3pt!-90:(Dp)$);
\draw[inkiron] ($(B)!0.51739253!(Dp)!3pt!90:(Dp)$) -- ($(B)!0.51739253!(Dp)!3pt!-90:(Dp)$);
\coordinate (E) at (-3.50000000,2.50000000);
\coordinate (F) at (1.51351351,-1.08108108);
\coordinate (M) at (-4.91891892,3.51351351);
\coordinate (N) at (-2.08108108,1.48648649);
\draw[indigo,dashed] (M) -- (F);
\draw[indigo,dashed] (C) -- (M);
\draw[indigo,dashed] (Dp) -- (N);
\pic[draw=selvage,angle radius=2.2mm] {right angle=C--M--B};
\pic[draw=selvage,angle radius=2.2mm] {right angle=Dp--N--B};
\pic[draw=selvage,angle radius=2.2mm] {right angle=B--F--D};
\fill[inkiron] (B) circle (1.05pt);
\node at (-0.46000000,-0.45000000) {$B$};
\fill[inkiron] (C) circle (1.05pt);
\node at (-6.45000000,2.10000000) {$C$};
\fill[inkiron] (Cp) circle (1.05pt);
\node at (-2.20000000,-6.42000000) {$C'$};
\fill[inkiron] (D) circle (1.05pt);
\node at (3.36000000,1.12000000) {$D$};
\fill[inkiron] (Dp) circle (1.05pt);
\node at (-0.68000000,3.26000000) {$D'$};
\fill[inkiron] (E) circle (1.05pt);
\node at (-3.46000000,2.90000000) {$E$};
\fill[inkiron] (F) circle (1.05pt);
\node at (1.81351351,-1.31108108) {$F$};
\fill[inkiron] (M) circle (1.05pt);
\node at (-5.01891892,3.91351351) {$M$};
\fill[inkiron] (N) circle (1.05pt);
\node at (-2.20108108,1.12648649) {$N$};
\end{tikzpicture}}


\begin{document}
\loomcover{双等腰模型}{婆罗摩笈多模型 · 单章核对稿}{Hence}{2026年10月}
\clearpage
{\small\tableofcontents}
\clearpage
\setcounter{section}{-1}
\section{一线三等角：快速回顾}
先看一幅熟悉的图：\(A\)、\(B\)、\(C\) 在同一直线上，\(B\) 在 \(A\)、\(C\) 之间；\(D\)、\(E\) 在直线 \(AC\) 的同侧，且
\[
\angle DAB=\angle DBE=\angle BCE=90^\circ,\qquad BD=BE.
\]
\diagram{\diReview}

直线上的三个角满足
\[
\angle DBA+\angle DBE+\angle EBC=180^\circ,
\]
所以 \(\angle DBA+\angle EBC=90^\circ\)。而在直角三角形 \(DAB\) 中，
\(\angle BDA+\angle DBA=90^\circ\)，故 \(\angle BDA=\angle EBC\)。

再结合 \(\angle DAB=\angle BCE\) 和 \(BD=BE\)，得到
\[
\triangle DBA\cong\triangle BEC\quad(\mathrm{AAS}),
\]
于是 \(\boldsymbol{AB=CE}\)、\(\boldsymbol{AD=BC}\)。注意对应顺序：
\(D\leftrightarrow B\)、\(B\leftrightarrow E\)、\(A\leftrightarrow C\)。

\begin{strand}
\textbf{识别顺序}\quad 一条直线 \(\longrightarrow\) 三个直角
\(\longrightarrow\) 等斜边 \(\longrightarrow\) 全等后交换对应直角边。

本章只回顾这一种直角全等型，下面“作双垂”的证明会连续用到它两次。
\end{strand}
\clearpage

\section{双等腰模型的三项结论}
\block{共用构型}
如图，\(\triangle CBC'\) 和 \(\triangle DBD'\) 都是以 \(B\) 为直角顶点的等腰直角三角形，即
\[
BC=BC',\quad BD=BD',\quad
\angle CBC'=\angle DBD'=90^\circ.
\]
各点方向按图取定：射线 \(BD\)、\(BD'\)、\(BC\)、\(BC'\) 绕 \(B\) 依次排列，\(\angle CBD'\) 为锐角。从 \(BC\) 到 \(BC'\)、从 \(BD\) 到 \(BD'\) 的 \(90^\circ\) 转动方向相同。以下三题共用这一构型，分别作辅助线。
\diagram{\diModel}

\begin{problem}{题 1 等面积从哪里来}
求证：\(S_{\triangle CBD'}=S_{\triangle C'BD}\)。
\hint{已经有 \(BC=BC'\)。分别以 \(BC\)、\(BC'\) 为底，试着比较两条高。}
\end{problem}

\begin{problem}{题 2 已知中点，证明垂直}
\(E\) 是 \(CD'\) 的中点。连接 \(BE\)，求证：\(BE\perp C'D\)。
\hint{中点条件能否通过倍长 \(BE\)，转化成一组八字形全等？}
\end{problem}

\begin{problem}{题 3 已知垂直，证明中点}
过 \(B\) 作 \(BF\perp C'D\)，\(F\) 为垂足；直线 \(BF\) 与 \(CD'\) 相交于 \(E\)。求证：\(CE=ED'\)。
\hint{还不知道 \(E\) 是中点。分别从 \(C\)、\(D'\) 向直线 \(BF\) 作垂线，看看会出现几组一线三等角。}
\end{problem}
\clearpage

\section{集中解析}
\begin{solution}{题 1 · 解析}
\think 两个三角形不一定全等，等面积也不必比较三边。先用已知的 \(BC=BC'\) 作等底，再把问题转成等高。

\solve 过 \(D'\) 作 \(D'M\perp BC\)，\(M\) 为垂足；过 \(D\) 作 \(DN\perp BC'\)，\(N\) 为垂足。本图中 \(M\) 在射线 \(BC\) 上，\(N\) 在 \(C'B\) 向 \(B\) 外的延长线上。
\diagram{\diArea}

由两个公共顶角都是直角，按图有
\[
\angle CBD'=180^\circ-\angle C'BD.
\]
由于 \(BM\) 与 \(BC\) 同向、\(BN\) 与 \(BC'\) 反向，
\[
\angle MBD'=\angle CBD',\qquad
\angle NBD=180^\circ-\angle C'BD,
\]
所以 \(\angle MBD'=\angle NBD\)。结合 \(\angle BMD'=\angle BND=90^\circ\)、\(BD'=BD\)，得到
\[
\triangle BMD'\cong\triangle BND\quad(\mathrm{AAS}),\qquad D'M=DN.
\]
从而
\[
S_{\triangle CBD'}=\frac12 BC\cdot D'M
=\frac12 BC'\cdot DN=S_{\triangle C'BD}.
\]

\insight{等面积先找等底；高不现成，就作垂线，把等高交给直角三角形全等。}
\end{solution}
\clearpage

\begin{solution}{题 2 · 解析}
\think \(E\) 是 \(CD'\) 的中点，但这个等长条件还不能和 \(BC=BC'\)、\(BD=BD'\) 放进同一组全等。倍长 \(BE\)，把 \(BD'\) 搬到 \(C\) 处，再比较一个大三角形。

\solve 延长 \(BE\) 到 \(G\)，使 \(EG=EB\)，连接 \(CG\)。于是 \(B\)、\(E\)、\(G\) 共线，\(CE=ED'\)，\(\angle CEG=\angle D'EB\)。
\diagram{\diMidpoint}

\solvehead{第一步：用八字形全等搬运一条边}
\[
\triangle CEG\cong\triangle D'EB\quad(\mathrm{SAS}).
\]
因此 \(CG=D'B=DB\)，且 \(\angle ECG=\angle ED'B\)，故 \(CG\parallel BD'\)。按图，射线 \(CG\) 与 \(BD'\) 同向。

\solvehead{第二步：比较一组大三角形}
由于 \(CB\) 与 \(BC\) 反向，\(CG\) 与 \(BD'\) 同向，
\[
\angle GCB=180^\circ-\angle CBD'=\angle DBC'.
\]
结合 \(CG=DB\)、\(CB=BC'\)，得到
\[
\triangle GCB\cong\triangle DBC'\quad(\mathrm{SAS}).
\]
于是 \(\angle GBC=\angle DC'B\)。

\solvehead{第三步：用对应角推出垂直}
延长 \(EB\) 过 \(B\)，与 \(C'D\) 相交于 \(H\)。按图，\(BH\)、\(BG\) 是反向射线，\(\angle CBC'=90^\circ\)，所以
\[
\angle C'BH+\angle GBC=90^\circ.
\]
又因 \(H\) 在 \(C'D\) 上，\(\angle BC'H=\angle DC'B=\angle GBC\)。在 \(\triangle BC'H\) 中，\(B\)、\(C'\) 两个角之和为 \(90^\circ\)，故 \(\angle BHC'=90^\circ\)，即 \(BE\perp C'D\)。

\solvehead{顺手得到的长度关系}
第二次全等还给出 \(GB=DC'\)。因为 \(GB=2BE\)，故
\[
\boxed{2BE=C'D}.
\]

\insight{倍长中线先搬运等长边，再用大三角形全等同时取得角度与长度关系。}
\end{solution}
\clearpage

\begin{solution}{题 3 · 解析}
\think 这一题和题 2 的已知、结论交换了位置。现在 \(CE=ED'\) 还是待证，不能把它放进倍长后的全等条件。题目给出了垂直，便从 \(C\)、\(D'\) 向 \(BF\) 作双垂，连续用两次章首模型。

\solve 过 \(C\) 作 \(CM\perp BF\)，过 \(D'\) 作 \(D'N\perp BF\)，\(M\)、\(N\) 为垂足。本题的 \(M\)、\(N\) 与题 1 的垂足不同。
\diagram{\diConverse}

\solvehead{第一步：用 \(BC=BC'\)，得到 \(CM=BF\)}
\(M\)、\(B\)、\(F\) 共线，\(\angle CMB=\angle CBC'=\angle BFC'=90^\circ\)。由直线上的角和，
\[
\angle CBM+\angle C'BF=90^\circ.
\]
而在直角三角形 \(C'BF\) 中，\(\angle BC'F+\angle C'BF=90^\circ\)，故 \(\angle CBM=\angle BC'F\)。再结合 \(BC=BC'\)，得到
\[
\triangle BCM\cong\triangle C'BF\quad(\mathrm{AAS}),\qquad CM=BF.
\]
这正是章首的一线三等角：等斜边推出对应直角边相等。

\solvehead{第二步：用 \(BD'=BD\)，得到 \(D'N=BF\)}
\(N\)、\(B\)、\(F\) 共线，\(\angle D'NB=\angle D'BD=\angle BFD=90^\circ\)。同样，
\begin{align*}
\angle D'BN+\angle DBF&=90^\circ,\\
\angle BDF+\angle DBF&=90^\circ,
\end{align*}
所以 \(\angle D'BN=\angle BDF\)。结合 \(BD'=BD\)，得到
\[
\triangle BD'N\cong\triangle DBF\quad(\mathrm{AAS}),\qquad D'N=BF.
\]
于是 \(CM=D'N\)。

\solvehead{第三步：接上八字形全等，得到中点}
按图，\(C\)、\(E\)、\(D'\) 共线，\(M\)、\(E\)、\(N\) 共线，\(\angle CME=\angle D'NE=90^\circ\)，\(\angle CEM=\angle D'EN\)。结合 \(CM=D'N\)，得到
\[
\triangle CME\cong\triangle D'NE\quad(\mathrm{AAS}),\qquad CE=ED'.
\]
所以 \(E\) 是 \(CD'\) 的中点。

若 \(M\)、\(N\) 恰好都与 \(E\) 重合，则 \(CM=D'N\) 直接给出 \(CE=ED'\)，不再使用退化三角形的全等。

\insight{要从垂直证明中点，先用双垂把同一条 \(BF\) 搬到两侧，再用八字形全等收尾。}
\end{solution}
\clearpage

\section{模型总结与迁移}
\ifteachernotes
\block{三个条件，三条路线}
\begin{center}\small
\renewcommand{\arraystretch}{1.3}
\begin{tabularx}{\linewidth}{@{}p{.26\linewidth}p{.22\linewidth}X@{}}
\toprule
已知／目标 & 辅助线 & 真正起作用的关系\\
\midrule
两等腰直角，证明等面积 & 以等长边作底，分别作高 & 等底＋直角三角形全等得到等高\\
\addlinespace
已知中点，证明垂直 & 倍长 \(BE\) & 八字形全等搬运边，再比较大三角形\\
\addlinespace
已知垂直，证明中点 & 从 \(C\)、\(D'\) 向 \(BE\) 作双垂 & 两次一线三等角，再接八字形全等\\
\bottomrule
\end{tabularx}
\end{center}

在本章构型下，可以把题 2、题 3 合起来记：
\[
E\in CD':\qquad CE=ED'\ \Longleftrightarrow\ BE\perp C'D.
\]
题 2 还告诉我们：当 \(E\) 是 \(CD'\) 的中点时，\textbf{\(BE\) 既垂直于 \(C'D\)，又等于 \(C'D\) 的一半}。

\block{再看一次等面积}
题 3 的两组全等给出 \(CM=D'N=BF\)。以 \(BE\) 为底，\(\triangle CBE\) 和 \(\triangle D'BE\) 的高都是 \(BF\)，所以
\begin{align*}
S_{\triangle CBD'}
&=S_{\triangle CBE}+S_{\triangle D'BE}\\
&=BE\cdot BF.
\end{align*}
而题 2 已经得到 \(2BE=C'D\)，故
\[
S_{\triangle CBD'}=\frac12 C'D\cdot BF=S_{\triangle C'BD}.
\]
这给出等面积的另一种看法：倍长中线得到的长度关系，和双垂得到的高度关系，最终在面积公式里相遇。

\block{易错提醒}
\begin{enumerate}
\item 两个等腰直角三角形的方向要按图取定。只记两组等长和两个直角，不看方位，不能直接套用本章结论。
\item 高是到\textbf{直线}的距离，垂足可能在边的延长线上；题 1 的 \(N\) 就在 \(BC'\) 的延长线上。
\item 题 3 中的中点是结论，不能预先使用 \(CE=ED'\) 来证明全等。
\item “两块面积始终相等”不表示“每块面积是一个固定值”；构型变化时，两块面积可以一起变化。
\item 三题分别作图。题 1、题 3 的 \(M\)、\(N\) 含义不同；题 2 的 \(G\) 是倍长点，\(H\) 是交点。
\end{enumerate}
\else
\block{学生自主总结}
在每道题旁记录：已知条件、第一条辅助线、第一次全等搬运了什么、第二次全等得到什么。再尝试写出中点、垂直与长度之间的关系。
\vspace{8\baselineskip}
\fi

\end{document}
